\documentclass{article}
\usepackage{amsmath}
\usepackage{color,pxfonts,fix-cm}
\usepackage{latexsym}
\usepackage[mathletters]{ucs}
\DeclareUnicodeCharacter{8594}{$\rightarrow$}
\DeclareUnicodeCharacter{46}{\textperiodcentered}
\DeclareUnicodeCharacter{58}{$\colon$}
\DeclareUnicodeCharacter{215}{$\times$}
\usepackage[T1]{fontenc}
\usepackage[utf8x]{inputenc}
\usepackage{pict2e}
\usepackage{wasysym}
\usepackage[english]{babel}
\usepackage{tikz}
\pagestyle{empty}
\usepackage[margin=0in,paperwidth=595pt,paperheight=841pt]{geometry}
\begin{document}
\definecolor{color_29791}{rgb}{0,0,0}
\begin{tikzpicture}[overlay]\path(0pt,0pt);\end{tikzpicture}
\begin{picture}(-5,0)(2.5,0)
\put(172.151,-106.071){\fontsize{20.6625}{1}\usefont{T1}{cmr}{m}{n}\selectfont\color{color_29791}TD Réseaux : Mo dèle OSI}
\put(238.449,-141.559){\fontsize{14.3462}{1}\usefont{T1}{cmr}{m}{n}\selectfont\color{color_29791}A del Hac hac he}
\put(41.69301,-189.4089){\fontsize{17.2154}{1}\usefont{T1}{cmr}{b}{n}\selectfont\color{color_29791}P artie I : Questions de cours}
\put(61.85301,-215.6909){\fontsize{11.9552}{1}\usefont{T1}{cmr}{m}{n}\selectfont\color{color_29791}1. Quelles son t les 7 couc hes du mo dèle OSI et leur rôle principal ?}
\put(61.85301,-233.1029){\fontsize{11.9552}{1}\usefont{T1}{cmr}{m}{n}\selectfont\color{color_29791}2. Quelle est la d ifférence en tre le mo dèle OSI et le mo dèle TCP/IP ?}
\put(61.85301,-250.5149){\fontsize{11.9552}{1}\usefont{T1}{cmr}{m}{n}\selectfont\color{color_29791}3. Expliquez les notions suiv an tes : proto cole, service, primitiv e, en tité.}
\put(61.85301,-267.9269){\fontsize{11.9552}{1}\usefont{T1}{cmr}{m}{n}\selectfont\color{color_29791}4. Donnez un exemple c oncret de proto cole p our c haque couc he.}
\put(61.85301,-285.3399){\fontsize{11.9552}{1}\usefont{T1}{cmr}{m}{n}\selectfont\color{color_29791}5. Qu’en tend-on par encapsulation et désencapsulation ?}
\put(41.69301,-324.6929){\fontsize{17.2154}{1}\usefont{T1}{cmr}{b}{n}\selectfont\color{color_29791}P artie I I : Exercices pratiques}
\put(41.69301,-354.4619){\fontsize{14.3462}{1}\usefont{T1}{cmr}{b}{n}\selectfont\color{color_29791}Exercice 1 : Couc he Ph ysique}
\put(61.85301,-376.6269){\fontsize{11.9552}{1}\usefont{T1}{cmr}{m}{n}\selectfont\color{color_29791}1. Un fic hier de 1 Go est transféré sur un lien de 1 Gb/s. Calculez le temps de transmission.}
\put(61.85301,-394.0389){\fontsize{11.9552}{1}\usefont{T1}{cmr}{m}{n}\selectfont\color{color_29791}2. Une trame de 1600 o ctets con ti e n t 1500 o ctets de données utiles. Calculez le p ourcen tage}
\put(76.80801,-408.4849){\fontsize{11.9552}{1}\usefont{T1}{cmr}{m}{n}\selectfont\color{color_29791}d’o v erhead.}
\put(61.85301,-425.8969){\fontsize{11.9552}{1}\usefont{T1}{cmr}{m}{n}\selectfont\color{color_29791}3.}
\put(407.8807,-425.8969){\fontsize{11.9552}{1}\usefont{T1}{cmr}{m}{n}\selectfont\color{color_29791}}
\put(411.7781,-425.8969){\fontsize{11.9552}{1}\usefont{T1}{cmr}{m}{n}\selectfont\color{color_29791}2}
\put(417.6314,-425.8969){\fontsize{11.9552}{1}\usefont{T1}{cmr}{m}{it}\selectfont\color{color_29791}}
\put(420.2855,-425.8969){\fontsize{11.9552}{1}\usefont{T1}{cmr}{m}{it}\selectfont\color{color_29791}×}
\put(429.5842,-425.8969){\fontsize{11.9552}{1}\usefont{T1}{cmr}{m}{n}\selectfont\color{color_29791}}
\put(432.2502,-425.8969){\fontsize{11.9552}{1}\usefont{T1}{cmr}{m}{n}\selectfont\color{color_29791}10}
\put(443.952,-421.5589){\fontsize{7.9701}{1}\usefont{T1}{cmr}{m}{n}\selectfont\color{color_29791}8}
\put(452.585,-425.8969){\fontsize{11.9552}{1}\usefont{T1}{cmr}{m}{n}\selectfont\color{color_29791}m/s).}
\put(61.85303,-443.3089){\fontsize{11.9552}{1}\usefont{T1}{cmr}{m}{n}\selectfont\color{color_29791}4. Sur 10 millions de bits transmis, 5000 son t erronés. Calculez le BER.}
\put(41.69303,-477.3919){\fontsize{14.3462}{1}\usefont{T1}{cmr}{b}{n}\selectfont\color{color_29791}Exercice 2 : Couc he Liaison de données}
\put(61.85303,-499.5569){\fontsize{11.9552}{1}\usefont{T1}{cmr}{m}{n}\selectfont\color{color_29791}1. Expliquer le rôle d e l’adresse MA C.}
\put(61.85303,-516.9689){\fontsize{11.9552}{1}\usefont{T1}{cmr}{m}{n}\selectfont\color{color_29791}2. Quelle est la d ifférence en tre un switc h et un p on t (bridge) ?}
\put(61.85303,-534.3809){\fontsize{11.9552}{1}\usefont{T1}{cmr}{m}{n}\selectfont\color{color_29791}3. Une trame Ethernet fait 1518 o ctets. Sac han t que l’en tête est de 18 o ctets, calculez le}
\put(76.80803,-548.8269){\fontsize{11.9552}{1}\usefont{T1}{cmr}{m}{n}\selectfont\color{color_29791}p ourcen tage de données utiles.}
\put(41.69303,-582.9089){\fontsize{14.3462}{1}\usefont{T1}{cmr}{b}{n}\selectfont\color{color_29791}Exercice 3 : Couc he Réseau}
\put(61.85303,-605.0739){\fontsize{11.9552}{1}\usefont{T1}{cmr}{m}{n}\selectfont\color{color_29791}1. Un message de 5000 o ctets doit être transmis en paquets IP de 1500 o ctets maxim um.}
\put(76.80803,-619.5199){\fontsize{11.9552}{1}\usefont{T1}{cmr}{m}{n}\selectfont\color{color_29791}Déterminez le nom bre de fragmen ts générés.}
\put(61.85303,-636.9319){\fontsize{11.9552}{1}\usefont{T1}{cmr}{m}{n}\selectfont\color{color_29791}2. Ca l c u lez le temps nécessaire p our transmettre un paquet de 1200 o ctets sur u n lien de}
\put(76.80803,-651.3779){\fontsize{11.9552}{1}\usefont{T1}{cmr}{m}{n}\selectfont\color{color_29791}2 Mb/s.}
\put(61.85303,-668.7899){\fontsize{11.9552}{1}\usefont{T1}{cmr}{m}{n}\selectfont\color{color_29791}3. Donnez deux différences en tre IPv4 et IPv6.}
\put(41.69303,-702.8719){\fontsize{14.3462}{1}\usefont{T1}{cmr}{b}{n}\selectfont\color{color_29791}Exercice 4 : Couc he T ransp ort}
\put(61.85303,-725.0369){\fontsize{11.9552}{1}\usefont{T1}{cmr}{m}{n}\selectfont\color{color_29791}1.}
\put(383.0128,-725.0369){\fontsize{11.9552}{1}\usefont{T1}{cmr}{m}{it}\selectfont\color{color_29791}}
\put(386.3722,-725.0369){\fontsize{11.9552}{1}\usefont{T1}{cmr}{m}{it}\selectfont\color{color_29791}E}
\put(395.036,-726.8299){\fontsize{7.9701}{1}\usefont{T1}{cmr}{m}{n}\selectfont\color{color_29791}1}
\put(399.768,-725.0369){\fontsize{11.9552}{1}\usefont{T1}{cmr}{m}{it}\selectfont\color{color_29791}..E}
\put(414.937,-726.8299){\fontsize{7.9701}{1}\usefont{T1}{cmr}{m}{n}\selectfont\color{color_29791}4}
\put(419.669,-725.0369){\fontsize{11.9552}{1}\usefont{T1}{cmr}{m}{n}\selectfont\color{color_29791},}
\put(485.5111,-725.0369){\fontsize{11.9552}{1}\usefont{T1}{cmr}{m}{it}\selectfont\color{color_29791}}
\put(488.8705,-725.0369){\fontsize{11.9552}{1}\usefont{T1}{cmr}{m}{it}\selectfont\color{color_29791}C}
\put(497.248,-726.8299){\fontsize{7.9701}{1}\usefont{T1}{cmr}{m}{n}\selectfont\color{color_29791}1}
\put(501.98,-725.0369){\fontsize{11.9552}{1}\usefont{T1}{cmr}{m}{it}\selectfont\color{color_29791}, C}
\put(515.6,-726.8299){\fontsize{7.9701}{1}\usefont{T1}{cmr}{m}{n}\selectfont\color{color_29791}0}
\put(520.332,-725.0369){\fontsize{11.9552}{1}\usefont{T1}{cmr}{m}{n}\selectfont\color{color_29791},}
\put(76.80804,-739.4829){\fontsize{11.9552}{1}\usefont{T1}{cmr}{m}{n}\selectfont\color{color_29791}sortie}
\put(104.831,-739.4829){\fontsize{11.9552}{1}\usefont{T1}{cmr}{m}{it}\selectfont\color{color_29791}}
\put(108.7284,-739.4829){\fontsize{11.9552}{1}\usefont{T1}{cmr}{m}{it}\selectfont\color{color_29791}S}
\put(115.9278,-739.4829){\fontsize{11.9552}{1}\usefont{T1}{cmr}{m}{n}\selectfont\color{color_29791}}
\put(116.6332,-739.4829){\fontsize{11.9552}{1}\usefont{T1}{cmr}{m}{n}\selectfont\color{color_29791}).}
\put(61.85304,-756.8949){\fontsize{11.9552}{1}\usefont{T1}{cmr}{m}{n}\selectfont\color{color_29791}2. Donnez l’équation logique corresp ondan te.}
\put(61.85304,-774.3069){\fontsize{11.9552}{1}\usefont{T1}{cmr}{m}{n}\selectfont\color{color_29791}3. Prop osez un exemple où l’UDP est préférable au TCP .}
\put(279.712,-804.1949){\fontsize{11.9552}{1}\usefont{T1}{cmr}{m}{n}\selectfont\color{color_29791}1}
\end{picture}
\newpage
\begin{tikzpicture}[overlay]\path(0pt,0pt);\end{tikzpicture}
\begin{picture}(-5,0)(2.5,0)
\put(41.693,-57.758){\fontsize{14.3462}{1}\usefont{T1}{cmr}{b}{n}\selectfont\color{color_29791}Exercice 5 : Couc he Session}
\put(61.853,-79.92297){\fontsize{11.9552}{1}\usefont{T1}{cmr}{m}{n}\selectfont\color{color_29791}1. Donnez deux rôles de la couc he Session.}
\put(61.853,-98.48596){\fontsize{11.9552}{1}\usefont{T1}{cmr}{m}{n}\selectfont\color{color_29791}2. Un fic hier de 2 50 Mo doit être transféré a v ec des p oin ts de reprise tous les 50 Mo. La}
\put(76.808,-112.9319){\fontsize{11.9552}{1}\usefont{T1}{cmr}{m}{n}\selectfont\color{color_29791}coupure survien t à 187 Mo. Com bien d e Mo doiv en t être retransmis ?}
\put(41.693,-147.5889){\fontsize{14.3462}{1}\usefont{T1}{cmr}{b}{n}\selectfont\color{color_29791}Exercice 6 : Couc he Présen tation}
\put(61.853,-169.7539){\fontsize{11.9552}{1}\usefont{T1}{cmr}{m}{n}\selectfont\color{color_29791}1. Citez deux fonctions de la couc he Présen tation.}
\put(61.853,-188.3169){\fontsize{11.9552}{1}\usefont{T1}{cmr}{m}{n}\selectfont\color{color_29791}2. Un fic hier de 40 Mo est compressé à 15 Mo. Calculez le taux de co mp re s sion.}
\put(61.853,-206.8799){\fontsize{11.9552}{1}\usefont{T1}{cmr}{m}{n}\selectfont\color{color_29791}3. Co mp a rez les temps de transmission sur un lien de 20 Mb/s p our la v ers io n compressée}
\put(76.808,-221.3249){\fontsize{11.9552}{1}\usefont{T1}{cmr}{m}{n}\selectfont\color{color_29791}et non compressée.}
\put(41.693,-255.9829){\fontsize{14.3462}{1}\usefont{T1}{cmr}{b}{n}\selectfont\color{color_29791}Exercice 7 : Couc he Applica tion}
\put(61.853,-278.1479){\fontsize{11.9552}{1}\usefont{T1}{cmr}{m}{n}\selectfont\color{color_29791}1. Asso ciez les proto co l e s suiv an ts à leur usage : HTTP , SMTP , DNS, FTP , DHCP .}
\put(61.853,-296.7109){\fontsize{11.9552}{1}\usefont{T1}{cmr}{m}{n}\selectfont\color{color_29791}2. Un fic hier de 10 Mo est téléc hargé via HTTP a v ec une latence de 50 ms par requête.}
\put(76.808,-311.1559){\fontsize{11.9552}{1}\usefont{T1}{cmr}{m}{n}\selectfont\color{color_29791}Calculez le temps total si 5 requêtes son t nécessaires (transmission incluse : 8 s).}
\put(279.712,-804.1949){\fontsize{11.9552}{1}\usefont{T1}{cmr}{m}{n}\selectfont\color{color_29791}2}
\end{picture}
\end{document}